% 基础宏包
\usepackage{xeCJKfntef, xpinyin}
\usepackage{graphicx}
\usepackage{zhlipsum}
\usepackage{tabularray}

% exam-zh 相关宏包（使用相对路径）
\usepackage{../exam-zh-choices}
\usepackage{../exam-zh-question}
\usepackage{../exam-zh-symbols}
\usepackage{../exam-zh-chinese-english}
\usepackage{../exam-zh-textfigure}
\usepackage{../exam-zh-math}

% examsetup 命令定义
\ExplSyntaxOn
\NewDocumentCommand \examsetup { m }
  { \keys_set:nn { exam-zh } {#1} }
\ExplSyntaxOff

% 图片路径
\graphicspath{{figures}{../doc/figures}}

% PDF 元数据
\hypersetup{
  pdftitle  = {exam-zh 入门指南：中国试卷 LaTeX 模板新手教程},
  pdfauthor = {夏康玮}
}

% 全角标点处理
\def\FSID{"{\xeCJKsetup{PunctStyle=banjiao}。}"} % U+3002
\def\FSFW{"{\xeCJKsetup{PunctStyle=banjiao}．}"} % U+FF0E
\def\COFW{"{\xeCJKsetup{PunctStyle=banjiao}：}"} % U+FF1A
\def\SCFW{"{\xeCJKsetup{PunctStyle=banjiao}；}"} % U+FF1B


% 文档标题和作者信息
\title{\textcolor{MaterialIndigo800}{%
  \textbf{exam-zh 入门指南\\[0.5em]
  \large 中国试卷 \LaTeX \xpinyin[font=\sffamily,format=\color{MaterialIndigo800}]{模}{mu2}板新手教程}}}
\author{%
  夏康玮\thanks{%
    李泽平构建了 \cls{exam-zh} 的最初的基本框架；张庭瑄开发 \cls{exam-zh-font} 模块；郭李军开发了连线题环境}
}
\date{\DocDate\quad \DocVersion%
  \thanks{%
    \textbf{入门文档}（推荐新手阅读）\\
    完整API参考见 \file{exam-zh-doc.pdf}\\[0.5em]
    GitHub: \url{https://github.com/xkwxdyy/exam-zh} \\
    Gitee（国内镜像）: \url{https://gitee.com/xkwxdyy/exam-zh} \\
    \hspace*{1.5em} QQ 用户交流群：652500180
  }
}

% 评分框命令（复用）
\ExplSyntaxOn
\NewDocumentCommand { \scoringbox } { s }
  {
    \IfBooleanTF {#1}
      { \__examzh_scoringbox_onecolumn: }
      { \__examzh_scoringbox_twocolumn: }
  }
\cs_new_protected:Nn \__examzh_scoringbox_twocolumn:
  {
    \begin{tabular}{|c|c|}
      \hline
      得分 & \rule{3em}{0pt}\rule[-0.7em]{0pt}{2em} \\\hline
      阅卷人 & \rule{3em}{0pt}\rule[-0.7em]{0pt}{2em} \\\hline
    \end{tabular}
  }
\cs_new_protected:Nn \__examzh_scoringbox_onecolumn:
  {
    \begin{tabular}{|c|}
      \hline
      得分\rule[-0.7em]{0pt}{2em} \\\hline
      \rule[-0.7em]{0pt}{2em} \\\hline
    \end{tabular}
  }
\ExplSyntaxOff

% 示例环境后重置题号
\AddToHook{env/latexexample/after}
  {%
    \examsetup{
      question/index=1
    }
  }
